\section{图像生成与编辑：从理解到创造}

\subsection{Create Image：接近真实的生成质量}

Qwen 于 2025 年开始图像生成方向，Create Image 系列在几个月内取得了显著进展：

\begin{itemize}
    \item 8 月版本：AI 感很重
    \item 12 月版本：接近照片级真实感（近乎``离谱''的逼真度）
    \item 内部 Arena 排名：超越最佳开源模型，接近 GPT-4.5 和 Gemini 3 Pro
    \item 特色能力：一次性生成带文字的分镜图（12 格漫画+对话文字，非拼接而是单张生成）
\end{itemize}

\subsection{图像编辑：比生成更刚需}

\begin{importantbox}{开源社区的意外发现}
Qwen 团队发现，图像\textbf{编辑}可能是比生成更大的需求。一个典型案例：用户想``放下图中人物的右手''，编辑后将两张图叠加时发现人物位置发生了偏移（shift）。这对 PS 专业用户来说是不可接受的精度问题。

2511 版本专门解决了这一问题——编辑后人物保持在原位，blending 后几乎无偏差。这个 case 来自开源社区的反馈，``如果不是开源社区告诉我们，可能这辈子都想不到有这个问题''。
\end{importantbox}

其他编辑能力包括：光线调整、物体移除、风格迁移、镜头旋转等。用户对图片物理合理性（如光影关系）的要求远超算法工程师的直觉。

\begin{knowledgebox}{理解-生成一体化的价值}
一个令人兴奋的应用场景是\textbf{辅助线绘制}：在教育领域教小朋友做几何题时，AI 光靠文字无法展示辅助线的画法——这需要图像生成能力与数学推理的结合。这是理解-生成一体化模型独特的价值所在。
\end{knowledgebox}

\subsection{本章小结}

图像生成和编辑是 Qwen ``全模态'' 愿景的重要拼图。开源社区的反馈在真实场景需求发现中起到了不可替代的作用。下一步是将编辑与生成统一到一个模型中。

